% Part I, list homomorphisms and the third homomorphism theorem: ~5 min

\begin{frame*}[vertical-alignment = top]
  \frametitle{List Homomorphisms}
  \vspace*{4mm}
  \hsblock{listhom.hs}
  \begin{center}
    \renewcommand\arraystretch{1.2}
    \begin{tabular}{llll}
      \code{h} & \code{h [x]} & \code{combine} & \code{h []} \\ \hline
      \code{sum} & \code{x} & \code{+} & \code{0} \\
      \code{maximum} & \code{x} & \code{max} & $-\infty$ \\
      \code{sort} & \code{[x]} & \code{merge} & \code{[]} \\
    \end{tabular}
  \end{center}
  \begin{onlyenv}<2->
    \textbf{First homomorphism theorem}: every such \code{h} is a
    \emph{map}, then a \emph{reduce} in any order. MapReduce.
  \end{onlyenv}
  \par
  \begin{onlyenv}<3->
    Any split gives the same answer, so \emph{how} to split is a
    \textbf{scheduling} decision (chunks per core, work stealing,
    cache-sized blocks), not a correctness one.
  \end{onlyenv}
  \par
  \refs<4>{R. Bird. An Introduction to the Theory of Lists. Logic of
    Programming and Calculi of Discrete Design, 1987.}
  \speakernote{Split anywhere, solve the pieces, combine \\
    sum, max, sort all qualify; first theorem (Bird): map then reduce \\
    Scheduling: the runtime picks the split for speed; correctness is
    settled once \\
    Click: reference}
\end{frame*}

\begin{frame*}
  \frametitle{Where Does the Combiner Come From?}
  \mainblock{digits.java}
  \begin{itemize}[<+->]
    \item Java's contract: the combiner is associative and agrees with the
      accumulator. Finding it is \textbf{your} job.
    \item Here it needed extra state, \code{scale}: a \textbf{lifting}.
  \end{itemize}
  \speakernote{Digits to a number: the obvious loop \\
    Parallel streams: you must supply the combiner \\
    A wrong combiner: wrong answers, only in parallel \\
    val alone is not enough: needs scale; same in Spark, OpenMP}
\end{frame*}

\begin{frame*}[vertical-alignment = top]
  \frametitle{The Third Homomorphism Theorem}
  \vspace*{4mm}
  \hsblock{third.hs}
  \begin{onlyenv}<2->
    \hsblock{third-sort.hs}
    Insertion sort, $O(n^2)$ $\Longrightarrow$ mergesort, $O(n \log n)$.
  \end{onlyenv}
  \par
  \refs<3>{J. Gibbons. The Third Homomorphism Theorem. Journal of
    Functional Programming, 1996.\\
    (Conjectured by R. Bird, proved by L. Meertens, 1989.)}
  \speakernote{foldl: the loop from the left; foldr: the same from the
    right \\
    stepL, stepR: one element at a time, need not be associative \\
    Both loops exist: then an associative combine exists, built from h and
    g (an example input for each result) \\
    Insertion sort both ways, g = id; calculate merge: mergesort for free \\
    Bird conjectured, Meertens proved it on a train; Gibbons published}
\end{frame*}

\begin{frame*}
  \frametitle{Homomorphisms Are Divide and Conquer}
  \hsblock{hom.hs}
  \speakernote{The divide step: split in the middle \\
    The conquer step: the combiner \\
    Balanced call tree, depth log n: parallel \\
    Divide by splitting off the first element instead: a loop}
\end{frame*}

\begin{frame}
  \frametitle{Same Combiner, Another Split}
  \begin{center}
    \begin{tikzpicture}[alt={Two call trees for the same combiner on four
        elements: a list-shaped tree of depth four, which splits off the
        first element, and a balanced tree of depth two, which splits in
        halves. Both visit x0, x1, x2, x3 in the same order.},
      every node/.style = {box, font = \small}]
      \begin{scope}[level distance = 8mm, sibling distance = 12mm]
        \node {$\oplus$}
          child {node {$x_0$}}
          child {node {$\oplus$}
            child {node {$x_1$}}
            child {node {$\oplus$}
              child {node {$x_2$}}
              child {node {$x_3$}}}};
      \end{scope}
      \node[draw = none, fill = none, note] at (0, 0.8)
        {split off the first: a loop};
      \node[draw = none, fill = none, font = \Large] at (35mm, -1.2)
        {$=$};
      \begin{scope}[xshift = 70mm, level distance = 10mm,
          level 1/.style = {sibling distance = 24mm},
          level 2/.style = {sibling distance = 12mm}]
        \node {$\oplus$}
          child {node {$\oplus$} child {node {$x_0$}} child {node {$x_1$}}}
          child {node {$\oplus$} child {node {$x_2$}} child {node {$x_3$}}};
      \end{scope}
      \node[draw = none, fill = none, note] at (70mm, 0.8)
        {split in halves: parallel};
    \end{tikzpicture}
  \end{center}
  \begin{itemize}[<+->]
    \item Same leaves, same order, associative $\oplus$: \textbf{same
      result}. We may \textbf{change how we divide}.
    \item The third theorem \textbf{constructs} such an $\oplus$ from two
      loops whose steps need not be associative.
  \end{itemize}
  \speakernote{Left: a loop, depth n. Right: depth log n \\
    Associativity of the combiner makes them equal \\
    The loop steps (accumulator, element) are not associative; the theorem
    builds the associative combine \\
    Remember this: changing the divide step comes back at the end}
\end{frame}

\begin{frame}
  \frametitle{Extensions}
  \begin{center}
    \small
    \renewcommand\arraystretch{1.25}
    \begin{tabular}{@{}l@{\quad}l@{\quad}l@{}}
      1987 & Bird & 1st and 2nd theorems: map--reduce; both directions \\
      1996 & Gibbons, JFP & 3rd theorem published; mergesort from
        insertion sort \\
      1996 & Gorlatch, Euro-Par & parallelising scans and homomorphisms \\
      1997 & Hu, Iwasaki, Takeichi, TOPLAS & combiners by adding extra
        state \\
      2007 & Morita et al., PLDI & combiners found automatically \\
      2009 & Morihata et al., POPL & the same for trees \\
      2011 & Mu, Morihata, ICFP & generalised and dualised \\
      2017--21 & Farzan, Nicolet, PLDI & synthesised for loops and nested
        loops \\
    \end{tabular}
  \end{center}
  \refs<2>{R. Bird. An Introduction to the Theory of Lists. 1987.\\
    J. Gibbons. The Third Homomorphism Theorem. JFP 1996.\\
    S. Gorlatch. Systematic Efficient Parallelization of Scan and Other
    List Homomorphisms. Euro-Par 1996.\\
    Z. Hu, H. Iwasaki, M. Takeichi. Formal Derivation of Efficient Parallel
    Programs by Construction of List Homomorphisms. TOPLAS 1997.\\
    K. Morita, A. Morihata, K. Matsuzaki, Z. Hu, M. Takeichi. Automatic
    Inversion Generates Divide-and-Conquer Parallel Programs. PLDI 2007.\\
    A. Morihata, K. Matsuzaki, Z. Hu, M. Takeichi. The Third Homomorphism
    Theorem on Trees. POPL 2009.\\
    S.-C. Mu, A. Morihata. Generalising and Dualising the Third
    List-Homomorphism Theorem. ICFP 2011.\\
    A. Farzan, V. Nicolet. Synthesis of Divide and Conquer Parallelism for
    Loops. PLDI 2017.\\
    A. Farzan, V. Nicolet. Modular Divide-and-Conquer Parallelization of
    Nested Loops. PLDI 2019.}
  \speakernote{From pen-and-paper theorems to tools \\
    Extra state = lifting, as in the Java example (tupling) \\
    Morita et al.: derive g by program inversion \\
    Farzan and Nicolet: program synthesis for loops; back in Part III \\
    Click: references}
\end{frame}
